\documentclass[11pt]{article}
\usepackage{multirow}

\begin{document}

\section{Probing results}

% The caption comes before the tabular here and swallows the label, which is
% the wiring style this fixture exists to pin down.
As shown in Table \ref{tab:probe}, finetuning helps on dev but transfers less
well to test. The second table has no label at all and is referenced only by
position, so extraction has to synthesise an identity for it.

\begin{table}[htbp]
\caption{Accuracy (\%) on the \textsc{NP/Z} probe set. {\bf Best} per column in
bold.\label{tab:probe}}
\centering
\begin{tabular}{lrrr}
\hline
Method & Dev & Test & $\Delta$ \\
\hline
Baseline & 71.2 & 69.8 & -- \\
+ finetune & \textbf{78.4} & 76.1 & +6.3 \\
+ finetune, 3 seeds & 78.1 (0.4) & 75.9 (0.6) & +6.1 \\
\hline
\end{tabular}
\end{table}

A second, unlabelled table groups the same numbers by construction.

\begin{table}[htbp]
\centering
\begin{tabular}{lrr}
\hline
Construction & GPT-2-large & Pythia-1.4B \\
\hline
\multicolumn{3}{l}{\emph{Ambiguous}} \\
NP/Z & 1.42 & 1.19 \\
NP/S & 0.88 & 0.79 \\
\hline
\multicolumn{3}{l}{\emph{Unambiguous}} \\
NP/Z & 0.31 & 0.28 \\
NP/S & 0.24 & 0.21 \\
\hline
\end{tabular}
\caption{Surprisal deltas grouped by construction. No label on this one.}
\end{table}

\end{document}
